%general context
As academics, we spend a lot of time setting in our rooms, not necessary worrying
about our surroundings, not a presence context, but in a breatheable, healthy
context. Many studies[LINK TIL DEM HER] show that spending time in negatively impacts
health and might even go as far as affecting the way we think, indirectly affecting
out minds. Many companies release IoT devices to monitor our bodily functions, suchs as
heart rate, stress levels and insulin, and others have even released air purifiers, to combat
this very problem. However, a common denominator for these solutions is that they are either
very expensive or tie the user down in complicated software eco-systems, and often the cases
is a combination of both.\\
With the increasing availability of low-cost embedded systems and wireless communication technologies,
the adoption of IoT solutions has accelerated the application in envionmental monitoring.\\
%why iot systems are useful
IoT based monitoring platforms can be implemented to use lightweight communication protocols 
and scaleable data-processing architectures. MQTT is one such protocol. Its low overhead and 
subscribe/publish communication model has made a widely adopted model for IoT applications. 
This makes it attractive when combining it with time-series databases suchs as InfluxDB and
visualization platorms such as grafana.\\

%project solution
This project presents a design and implementation of a IoT-enabled envionmental monitoring 
and fan-control system based on a newer PIC16 microcontroller and a Rasbperry PI gateway.
The embedded subsystem aquires air quiality measurements from an ENS160 gas sensor, while 
simultaneously monitoring multiple PWM-controlled fans. Telemetry is transmitted through
a custom UART communications protocol to the Raspbery Pi, where data is published to MQTT, 
processing using Telegraf, stored in InfluxDb and visualized in Grafana dashboards.

%paper overview
The purpose of this paper is to demonstrate a low-cost, lightweight, embedded platform,
that can be integrated into larger, modular IoT architectures, for monitoring and Telemetry
visualization. The remainder of this reports describes the theoretical background, system
architecture, implementation details, experimental verification and future improvements.